% Main blueprint content. Mirrors §1--§3 of
%   R. Valencia, "Gamma Operators for Pochhammer Kernels"
% and tags every statement with its Lean counterpart so plasTeX can
% draw the formalization dependency graph.

\chapter*{Overview}

This blueprint tracks the formalization, in Lean~4 + mathlib,
of the paper
\emph{Gamma Operators for Pochhammer Kernels}.
The Lean development lives in the
\href{https://github.com/rvazdev-ex/gamma-pochhammer}{GammaPochhammer}
repository; this document mirrors its mathematical content statement
by statement.

The build state, as of the current commit, is:

\begin{itemize}
  \item Theorem~1 (determinant polynomial preserves negative-rooted) is
        proved modulo the two Pólya--Schur axioms (Lemmas~1 and 2).
  \item Theorem~2 (ladder formula) is proved unconditionally.
  \item Theorem~3 (centered-kernel classification) is proved modulo
        the Schur--Maló axiom (Lemma~2).
  \item Lemma~3 (the $\gamma$-representation of palindromic
        negative-rooted polynomials) is proved from first principles.
\end{itemize}

Every \emph{green} node in the dependency graph is fully formalized in
Lean.  Two red nodes remain: the two Pólya--frequency / Schur
classical results that mathlib does not yet have, namely Lemmas~1 and 2.

\chapter{Preliminaries}
\label{ch:preliminaries}

\section{Polynomials with real, non-positive zeros}

\begin{definition}[Real non-positive zeros]
  \label{def:HasOnlyRealNonposZeros}
  \lean{GammaPochhammer.HasOnlyRealNonposZeros}
  \leanok
  A real polynomial $p \in \RR[X]$ is said to have \emph{only real,
  non-positive zeros} when every complex root of $p$, regarded as a
  polynomial in $\CC[X]$, lies on the non-positive real axis
  $\RR_{\leq 0}$.
\end{definition}

\begin{definition}[Real strictly negative zeros]
  \label{def:HasOnlyRealNegativeZeros}
  \lean{GammaPochhammer.HasOnlyRealNegativeZeros}
  \leanok
  A real polynomial $p \in \RR[X]$ has \emph{only real, strictly
  negative zeros} when every complex root lies in $(-\infty, 0)$.
  Equivalently: $p(0) \neq 0$ and $p$ has only real, non-positive zeros.
\end{definition}

\begin{definition}[Degree marker]
  \label{def:HasDegree}
  \lean{GammaPochhammer.HasDegree}
  \leanok
  We write $\operatorname{HasDegree}(p, n)$ for the proposition that
  $p \in \RR[X]$ has natural degree exactly $n$ and is non-zero.
\end{definition}

\section{Pochhammer symbols and the determinant polynomial}

\begin{definition}[Pochhammer symbol]
  \label{def:pochhammer}
  \lean{GammaPochhammer.pochhammer}
  \leanok
  For $a \in \RR$ and $k \in \NN$, the rising factorial is
  \[
    \Poch{a}{k} \;=\; a(a+1)(a+2)\cdots(a+k-1)
        \;=\; \prod_{j=0}^{k-1}(a+j) \in \RR[X].
  \]
\end{definition}

\begin{definition}[Linear Pochhammer]
  \label{def:linearPochhammer}
  \lean{GammaPochhammer.linearPochhammer}
  \leanok
  $\operatorname{linearPochhammer}(b, a, k)$ is the polynomial obtained
  by substituting $bX + a$ into $\Poch{X}{k}$; used to model the
  argument of a rung kernel.
\end{definition}

\begin{definition}[Ordinary generating polynomial]
  \label{def:ordinaryGen}
  \lean{GammaPochhammer.ordinaryGen}
  \leanok
  For $n \in \NN$ and $f : \NN \to \RR$ we set
  \[
    \operatorname{ordinaryGen}(n, f) \;=\; \sum_{k=0}^{n} f(k)\, X^{k}.
  \]
\end{definition}

\begin{definition}[Hadamard product]
  \label{def:hadamard}
  \lean{GammaPochhammer.hadamard}
  \leanok
  The (truncated) Hadamard product of $A, B \in \RR[X]$ at order $n$ is
  \[
    \operatorname{hadamard}(n, A, B) \;=\; \sum_{k=0}^{n}
        a_k\, b_k\, X^{k}, \qquad
    a_k = [X^k] A,\ b_k = [X^k] B.
  \]
\end{definition}

\begin{definition}[Determinant kernel]
  \label{def:determinantKernel}
  \lean{GammaPochhammer.determinantKernel}
  \leanok
  $\Kdet{n}{k}(X) = \Poch{X}{k}\, \Poch{X}{n-k}$.
\end{definition}

\begin{definition}[Determinant polynomial $\Pn$]
  \label{def:determinantPolynomial}
  \lean{GammaPochhammer.determinantPolynomial}
  \leanok
  Given $f : \NN \to \RR$ and $n \in \NN$,
  \[
    \Pn(X; f) \;=\; \sum_{k=0}^{n} \binom{n}{k} f_k\, f_{n-k}\, \Kdet{n}{k}(X).
  \]
\end{definition}

\section{Palindromic polynomials and the \texorpdfstring{$\gamma$}{gamma}-basis}

\begin{definition}[Palindromic of degree $n$]
  \label{def:PalindromicOfDegree}
  \lean{GammaPochhammer.PalindromicOfDegree}
  \leanok
  $Q \in \RR[X]$ is \emph{palindromic of degree $n$} when it has natural
  degree $n$ and $[X^{n-k}]Q = [X^k]Q$ for all $0 \leq k \leq n$.
\end{definition}

\begin{definition}[$\gamma$-basis]
  \label{def:gammaBasis}
  \lean{GammaPochhammer.gammaBasis}
  \leanok
  For $n, j \in \NN$ we set $\gammabasis{n}{j} = X^{j}(1+X)^{\,n-2j}$.
  These are the building blocks of the $\gamma$-representation of any
  palindromic polynomial of degree~$n$ with only real, non-positive
  zeros.
\end{definition}

\begin{definition}[$\gamma$-expansion]
  \label{def:IsGammaExpansion}
  \lean{GammaPochhammer.IsGammaExpansion}
  \leanok
  $\operatorname{IsGammaExpansion}(n, m, Q, \gamma)$ asserts that
  \(
    Q \;=\; \sum_{j=0}^{m} \gamma_j\, \gammabasis{n}{j}.
  \)
\end{definition}

\begin{definition}[$\gamma$-polynomial]
  \label{def:gammaPolynomial}
  \lean{GammaPochhammer.gammaPolynomial}
  \leanok
  $\Gpoly{m}{\gamma} = \sum_{j=0}^{m} \gamma_j\, \gammabasis{n}{j}$,
  packaged as a Lean polynomial.
\end{definition}

\chapter{Theorem~1: the determinant polynomial}
\label{ch:theorem1}

\section{The two Pólya--Schur axioms}

These are the only two statements in the development that are still
axioms.  Both are classical Pólya--frequency results that mathlib does
not currently provide.

\begin{lemma}[Hadamard closure / Aissen--Schoenberg--Whitney]
  \label{lem:hadamard_closure}
  \lean{GammaPochhammer.hadamard_closure_for_negative_rooted}
  \leanok
  Let $A, B \in \RR[X]$ both have only real, strictly negative zeros and
  positive leading coefficient.  Then the truncated Hadamard product
  $\operatorname{hadamard}(n, A, B)$ has only real, strictly negative
  zeros.
\end{lemma}

The Lean statement is recorded as the axiom
\lean{GammaPochhammer.hadamard_closure_for_negative_rooted}
in \texttt{Basic.lean}; an axiom-free proof would require the
Aissen--Schoenberg--Whitney theory of totally-positive sequences.

\begin{lemma}[Schur--Maló: $\Lop$ preserves non-positive-rooted]
  \label{lem:schur_preserves_nonpos}
  \lean{GammaPochhammer.schur_preserves_nonpos}
  \leanok
  Let $B \in \RR[X]$ have only real, non-positive zeros.  Then for all
  $\ell \in \NN$ and $\alpha \geq 0$, the Schur transform
  $\Lop_{\ell,\alpha}(B)$ has only real, non-positive zeros.
\end{lemma}

Recorded in Lean as the axiom
\lean{GammaPochhammer.schur_preserves_nonpos}.  An axiom-free proof
would require Hermite--Kakeya--Obreschkoff (HKO) interlacing.

\subsection*{Mathlib-facing wrappers}

\begin{lemma}[Mathlib wrapper for Hadamard closure]
  \label{lem:Hadamard_closure_wrapper}
  \lean{GammaPochhammer.Hadamard_closure_for_negative_rooted_polynomials}
  \leanok
  \uses{lem:hadamard_closure}
  Convenience restatement of Lemma~\ref{lem:hadamard_closure} in the
  exact form consumed by the proof of Theorem~1.
\end{lemma}

\begin{lemma}[Mathlib wrapper for Schur--Maló]
  \label{lem:schur_transform_preserves_nonpos}
  \lean{GammaPochhammer.schur_transform_preserves_nonpos}
  \leanok
  \uses{lem:schur_preserves_nonpos}
  Convenience restatement of Lemma~\ref{lem:schur_preserves_nonpos}
  for the form used inside the proof of Theorem~1.
\end{lemma}

\section{Lemma~3: the \texorpdfstring{$\gamma$}{gamma}-representation}

\begin{lemma}[$\gamma$-representation]
  \label{lem:gamma_representation}
  \lean{GammaPochhammer.gamma_representation}
  \leanok
  \uses{def:PalindromicOfDegree, def:gammaBasis,
        def:IsGammaExpansion, def:HasOnlyRealNegativeZeros}
  Let $Q$ be palindromic of degree $n = 2m + \varepsilon$
  ($\varepsilon \in \{0,1\}$) with only real, non-positive zeros and
  positive leading coefficient.  Then there exist non-negative
  coefficients $\gamma_0, \ldots, \gamma_m \geq 0$ such that
  \(
    Q \;=\; \sum_{j=0}^{m} \gamma_j\, \gammabasis{n}{j}.
  \)
\end{lemma}

\begin{proof}
  \leanok
  \uses{def:PalindromicOfDegree, def:gammaBasis, def:HasOnlyRealNegativeZeros}
  Strong induction on the degree~$n$, using the paper's reciprocal-pair
  factorization.  The Lean proof is in
  \texttt{GammaPochhammer/GammaRep.lean} and uses no project-specific
  axioms.
\end{proof}

\begin{lemma}[$\gamma$-representation, negative-rooted form]
  \label{lem:gamma_representation_of_palindromic_negative_rooted}
  \lean{GammaPochhammer.gamma_representation_of_palindromic_negative_rooted}
  \leanok
  \uses{lem:gamma_representation, def:HasOnlyRealNegativeZeros}
  Strengthening of Lemma~\ref{lem:gamma_representation} where $Q$ has
  only strictly negative zeros; the same $\gamma$-expansion exists.
\end{lemma}

\section{Theorem~1}

\begin{theorem}[Theorem~1 --- $\Pn$ preserves real-negative-rooted families]
  \label{thm:main}
  \lean{GammaPochhammer.main_theorem}
  \leanok
  \uses{def:determinantPolynomial, def:HasOnlyRealNegativeZeros,
        lem:hadamard_closure, lem:schur_preserves_nonpos,
        lem:gamma_representation}
  Let $n \geq 2$ and $f : \NN \to \RR$ be such that the generating
  polynomial $\operatorname{ordinaryGen}(n,f)$ has only real, strictly
  negative zeros, positive leading coefficient, and $f_0 > 0$.  Then
  $\Pn(X; f)$ has only real, strictly negative zeros.
\end{theorem}

\begin{proof}
  \leanok
  \uses{def:determinantPolynomial, lem:Hadamard_closure_wrapper,
        lem:schur_transform_preserves_nonpos,
        lem:gamma_representation,
        lem:convolutionPoly_palindromic,
        lem:convolutionPoly_hasOnlyRealNegativeZeros,
        lem:determinantPolynomial_eq_rungOperator_convolutionPoly,
        lem:cauchy_schwarz_f1_fn1_gt_f0_fn,
        lem:rung_action,
        thm:ladder_formula}
  Follow the paper's proof of Theorem~1:
  \begin{enumerate}
    \item Form the convolution polynomial
          $Q = \sum_k \binom{n}{k} f_k f_{n-k} X^k$
          (Lemma~\ref{lem:convolutionPoly_palindromic}).
    \item Show $Q$ is palindromic of degree~$n$ with only real,
          strictly negative zeros
          (Lemma~\ref{lem:convolutionPoly_hasOnlyRealNegativeZeros}),
          using Hadamard closure (Lemma~\ref{lem:hadamard_closure})
          twice and the
          $(1+X)^n$ baseline.
    \item Identify $\Pn$ with the $s=1$, $\mu=0$ rung operator applied
          to $Q$ (Lemma~\ref{lem:determinantPolynomial_eq_rungOperator_convolutionPoly}).
    \item Apply the rung-action formula
          (Lemma~\ref{lem:rung_action})
          together with the $\gamma$-representation
          (Lemma~\ref{lem:gamma_representation}) to express $\Pn$
          as a Schur transform of $Q$.
    \item Conclude by Lemma~\ref{lem:schur_preserves_nonpos}, using
          $f_1 f_{n-1} > f_0 f_n$
          (Lemma~\ref{lem:cauchy_schwarz_f1_fn1_gt_f0_fn})
          to exclude a zero at the origin.
  \end{enumerate}
  Full Lean proof in
  \texttt{GammaPochhammer/Determinant.lean}.
\end{proof}

\subsection{Vieta + convolution infrastructure}

These are paper-internal results used by the proof of Theorem~1.  All
are formalized in Lean and depend on no project-specific axiom.

\begin{lemma}[Convolution polynomial]
  \label{lem:convolutionPoly}
  \lean{GammaPochhammer.convolutionPoly}
  \leanok
  $\operatorname{convolutionPoly}(n,f)
   = \sum_{k=0}^{n} \binom{n}{k} f_k f_{n-k} X^k.$
\end{lemma}

\begin{lemma}[Convolution polynomial is palindromic]
  \label{lem:convolutionPoly_palindromic}
  \lean{GammaPochhammer.convolutionPoly_palindromic}
  \leanok
  \uses{lem:convolutionPoly, def:PalindromicOfDegree}
  $\operatorname{convolutionPoly}(n,f)$ is palindromic of degree~$n$.
\end{lemma}

\begin{lemma}[Convolution polynomial has only real negative zeros]
  \label{lem:convolutionPoly_hasOnlyRealNegativeZeros}
  \lean{GammaPochhammer.convolutionPoly_hasOnlyRealNegativeZeros}
  \leanok
  \uses{lem:convolutionPoly, lem:hadamard_closure,
        def:HasOnlyRealNegativeZeros}
  Under the hypotheses of Theorem~\ref{thm:main},
  $\operatorname{convolutionPoly}(n,f)$ has only real, strictly negative
  zeros.  Proved by two applications of Hadamard closure (with
  $\operatorname{ordinaryGen}$, its coefficient-reversal, and the
  baseline $(1+X)^n$).
\end{lemma}

\begin{lemma}[$\Pn$ as a rung operator]
  \label{lem:determinantPolynomial_eq_rungOperator_convolutionPoly}
  \lean{GammaPochhammer.determinantPolynomial_eq_rungOperator_convolutionPoly}
  \leanok
  \uses{def:determinantPolynomial, def:rungOperator, lem:convolutionPoly}
  $\Pn(X; f) = T_{1,0}^{(n)}\!\bigl(\operatorname{convolutionPoly}(n,f)\bigr)$.
\end{lemma}

\begin{lemma}[$\Pn(0)$]
  \label{lem:determinantPolynomial_eval_zero}
  \lean{GammaPochhammer.determinantPolynomial_eval_zero}
  \leanok
  \uses{def:determinantPolynomial}
  $\Pn(0; f) = n!\,(f_1 f_{n-1} - f_0 f_n)$.
\end{lemma}

\begin{lemma}[$f_1 f_{n-1} > f_0 f_n$]
  \label{lem:cauchy_schwarz_f1_fn1_gt_f0_fn}
  \lean{GammaPochhammer.cauchy_schwarz_f1_fn1_gt_f0_fn}
  \leanok
  \uses{def:ordinaryGen, def:HasOnlyRealNegativeZeros,
        lem:determinantPolynomial_eval_zero}
  Under the hypotheses of Theorem~\ref{thm:main} we have
  $f_1 f_{n-1} > f_0 f_n$.  This is the paper's
  ``$e_1 e_{n-1} > e_n$'' Newton-type inequality and prevents
  $\Pn$ from acquiring a root at the origin.
\end{lemma}

\begin{lemma}[Internal: derived Theorem~1]
  \label{lem:determinant_main_preserves_strict_proved}
  \lean{GammaPochhammer.determinant_main_preserves_strict_proved}
  \leanok
  \uses{lem:Hadamard_closure_wrapper,
        lem:schur_transform_preserves_nonpos,
        lem:gamma_representation,
        lem:convolutionPoly_hasOnlyRealNegativeZeros,
        lem:determinantPolynomial_eq_rungOperator_convolutionPoly,
        lem:cauchy_schwarz_f1_fn1_gt_f0_fn}
  Internal name for Theorem~\ref{thm:main} (the version that takes the
  full Lemma~1 and Lemma~2 hypotheses).
  \texttt{main\_theorem} is a thin user-facing wrapper.
\end{lemma}

\chapter{Theorem~2: the Pochhammer kernel ladder}
\label{ch:theorem2}

\section{Rung kernel and rung operator}

\begin{definition}[$\lambda_{s,q}$ coefficients]
  \label{def:lambda}
  \lean{GammaPochhammer.lambda}
  \leanok
  $\lcoef{s}{q} \in \RR$ is the alternating central-difference
  coefficient
  $\lcoef{s}{q} = (-1)^{q}\binom{2s}{s-q}$ (paper, eq.~$\lambda$).
\end{definition}

\begin{definition}[Rung kernel $K^{[s,\mu]}_{n,k}$]
  \label{def:rungKernel}
  \lean{GammaPochhammer.rungKernel}
  \leanok
  $\Krung{n}{k}{s}{\mu}(X)
   = \Poch{X+\mu}{k+s}\,\Poch{X+\mu}{n-k+s}$.
\end{definition}

\begin{definition}[Rung operator $T_{s,\mu}^{(n)}$]
  \label{def:rungOperator}
  \lean{GammaPochhammer.rungOperator}
  \leanok
  For $Q \in \RR[X]$, $T_{s,\mu}^{(n)}(Q) = \sum_{k=0}^{n}
        [X^k]Q \cdot \Krung{n}{k}{s}{\mu}.$
\end{definition}

\section{The $A380113$ triangle}

\begin{definition}[$A_{s,q}$ triangle]
  \label{def:A380113}
  \lean{GammaPochhammer.A380113}
  \leanok
  $\Atriangle{s}{q}$ is the OEIS sequence A380113, equal to
  $\binom{2s}{s-q}$ in the closed form below.
\end{definition}

\begin{lemma}[A380113 base case]
  \label{lem:A380113_zero_zero}
  \lean{GammaPochhammer.A380113_zero_zero}
  \leanok
  $\Atriangle{0}{0} = 1$.
\end{lemma}

\begin{lemma}[A380113 left edge]
  \label{lem:A380113_left_edge}
  \lean{GammaPochhammer.A380113_left_edge}
  \leanok
  For $s \neq 0$, $\Atriangle{s}{0}$ is given by the paper's left-edge
  formula.
\end{lemma}

\begin{lemma}[A380113 internal recursion]
  \label{lem:A380113_internal}
  \lean{GammaPochhammer.A380113_internal}
  \leanok
  For $q \neq 0$, $\Atriangle{s}{q}$ satisfies the paper's interior
  recurrence.
\end{lemma}

\section{Lemma~4 and Lemma~5}

\begin{lemma}[Lemma~4 --- central-difference identity]
  \label{lem:central_difference}
  \lean{GammaPochhammer.central_difference}
  \leanok
  \uses{def:lambda, def:pochhammer}
  The alternating sum
  $\sum_{q} \lcoef{s}{q} \Poch{w+q}{?}$ collapses to the paper's
  closed-form central-difference identity.
\end{lemma}

\begin{lemma}[Pointwise version]
  \label{lem:central_difference_identity}
  \lean{GammaPochhammer.central_difference_identity}
  \leanok
  \uses{lem:central_difference}
  Pointwise restatement of Lemma~4 used in the proof of Theorem~2.
\end{lemma}

\begin{lemma}[Lemma~5 --- rung action on the $\gamma$-basis]
  \label{lem:rung_action_on_gamma_basis}
  \lean{GammaPochhammer.rung_action_on_gamma_basis}
  \leanok
  \uses{def:rungOperator, def:gammaBasis, lem:central_difference}
  The rung operator $T_{s,\mu}^{(n)}$ acts on a single $\gamma$-basis
  element $\gammabasis{n}{j}$ by the explicit formula recorded in the
  paper (eq.~$T B_{n,j}$).
\end{lemma}

\begin{lemma}[Lemma~5 (rung action on a $\gamma$-expansion)]
  \label{lem:rung_action}
  \lean{GammaPochhammer.rung_action}
  \leanok
  \uses{lem:rung_action_on_gamma_basis, def:IsGammaExpansion}
  By linearity from Lemma~\ref{lem:rung_action_on_gamma_basis}.
\end{lemma}

\section{Theorem~2 and its corollaries}

\begin{theorem}[Theorem~2 --- ladder formula]
  \label{thm:ladder_formula}
  \lean{GammaPochhammer.ladder_formula}
  \leanok
  \uses{lem:rung_action, def:rungOperator, def:gammaBasis}
  Statement of the paper's ladder formula: applied to a $\gamma$-expansion,
  the rung operator $T_{s,\mu}^{(n)}$ produces another rung kernel with
  shifted parameters.  (Full statement in the source.)
\end{theorem}

\begin{theorem}[Pochhammer kernel ladder, full form]
  \label{thm:pochhammer_kernel_ladder_formula}
  \lean{GammaPochhammer.pochhammer_kernel_ladder_formula}
  \leanok
  \uses{thm:ladder_formula}
  The fully-expanded form of Theorem~\ref{thm:ladder_formula}.
\end{theorem}

\begin{theorem}[Kernel ladder, downstream]
  \label{thm:pochhammer_kernel_ladder}
  \lean{GammaPochhammer.pochhammer_kernel_ladder}
  \leanok
  \uses{thm:ladder_formula, lem:gamma_representation}
  Downstream form of Theorem~2 used by Corollary~1 and Corollary~2.
\end{theorem}

\begin{theorem}[Strict-shift form]
  \label{thm:pochhammer_kernel_ladder_strict_shift}
  \lean{GammaPochhammer.pochhammer_kernel_ladder_strict_shift}
  \leanok
  \uses{thm:pochhammer_kernel_ladder}
  Variant in which the shift parameter is strict.
\end{theorem}

\begin{corollary}[Corollary~1 --- original kernel is palindromic]
  \label{cor:original_kernel_palindromic}
  \lean{GammaPochhammer.original_kernel_palindromic}
  \leanok
  \uses{thm:pochhammer_kernel_ladder, def:PalindromicOfDegree}
  For all admissible parameters, the original Pochhammer kernel is
  palindromic of the appropriate degree.
\end{corollary}

\begin{corollary}[Corollary~2 --- single-product kernel at $s=0$]
  \label{cor:single_product_kernel}
  \lean{GammaPochhammer.single_product_kernel}
  \leanok
  \uses{thm:pochhammer_kernel_ladder}
  At $s = 0$ the rung kernel collapses to a single Pochhammer product.
\end{corollary}

\chapter{Theorem~3: classification of universally-admissible centered kernels}
\label{ch:theorem3}

\section{The centered two-product family}

\begin{definition}[Centered kernel $K^{a,c}_{n,k}$]
  \label{def:centeredKernel}
  \lean{GammaPochhammer.centeredKernel}
  \leanok
  $\Kcent{n}{k}{a}{c}(X) = \Poch{X+a}{k}\,\Poch{X+c}{n-k}$.
\end{definition}

\begin{definition}[Centered operator]
  \label{def:centeredOperator}
  \lean{GammaPochhammer.centeredOperator}
  \leanok
  $T^{a,c}_n(Q) = \sum_{k=0}^{n} [X^k] Q \cdot \Kcent{n}{k}{a}{c}$.
\end{definition}

\begin{definition}[Universally admissible]
  \label{def:UniversallyAdmissibleCentered}
  \lean{GammaPochhammer.UniversallyAdmissibleCentered}
  \leanok
  A pair $(a, c) \in \RR^2$ is \emph{universally admissible} when
  $T^{a,c}_n$ sends every palindromic real-rooted $Q$ to a polynomial
  with only real zeros.  (Full Lean definition in
  \texttt{Basic.lean}.)
\end{definition}

\section{Theorem~3}

\begin{theorem}[Theorem~3 --- centered classification]
  \label{thm:centered_classification}
  \lean{GammaPochhammer.centered_balanced_classification}
  \leanok
  \uses{def:UniversallyAdmissibleCentered, def:centeredKernel,
        lem:schur_preserves_nonpos, cor:original_kernel_palindromic}
  Let $a, c \in \RR$ with $a^2 \neq c^2$.  Then $(a,c)$ is universally
  admissible if and only if $\{a^2, c^2\} = \{0, 1\}$ as subsets of
  $\RR$.
\end{theorem}

\begin{proof}
  \leanok
  \uses{lem:schur_preserves_nonpos, cor:original_kernel_palindromic,
        lem:gamma_representation,
        thm:ladder_formula}
  Follows the paper's discriminant + intermediate-value-theorem
  argument: $\sigma = u + v = 1$ is forced via two $\lambda$-tests on a
  perturbation $Q_n^{(4)}(\lambda) = (\text{quadFactor}\,\lambda)^2$;
  $u = 0$ is then forced by an $\varepsilon$-test on
  $Q_n^{(6)}(\varepsilon) = (\text{quadFactor}\,\varepsilon)^3$ and
  the intermediate value theorem.  Sufficiency reduces to the
  determinant-kernel case using sign and palindromic symmetries via
  Corollary~\ref{cor:original_kernel_palindromic}.  Full Lean proof in
  \texttt{GammaPochhammer/Classification.lean}.
\end{proof}

\begin{lemma}[Internal name]
  \label{lem:centered_classification_proved}
  \lean{GammaPochhammer.centered_classification_proved}
  \leanok
  \uses{thm:centered_classification}
  Internal name for Theorem~\ref{thm:centered_classification}.
  \texttt{centered\_balanced\_classification} is the user-facing
  alias.
\end{lemma}
